\section{Bootstrapping the module loader}
\label{module-loader}

% PAPER-REVIEW[codex:PR-STRUCTURE-0002,claude:PR-TECHNICAL-0023] pass=structure,technical state=open
% Issue (codex): The bridge from individually bootstrapped macros to the loader's complete environment is missing.
% Evidence (codex): The introduction says let* details are omitted although let* was just defined; subsequent code also uses list, concat, sym-item-pair, !and and kv-fetch without a dependency inventory.
% Proposal (codex): Add a compact dependency map that separates already-defined utilities, omitted routine helpers and the order in which their scopes are built; explain the input/output contracts of omitted helpers, especially sym-item-pair.
% Reason (codex): This makes the top-down presentation navigable and separates conceptual dependencies from implementation omissions.
% Issue (claude): [Missing information] The central artefact cannot be reproduced from the paper.
% Evidence (claude): module-loader, process-import-items, kv-fetch, sym-item-pair, concat and the bootstrap files
%   core/bootstrap/*.ild are omitted, and no implementation, repository or artifact is referenced, although a naive
%   interpreter is mentioned (Section 7.2.1). The body ends on p.10 of 12 allowed pages (R1), and appendices are unlimited.
% Proposal (claude): Add a repository/artifact link (allowed under single-blind review, R3) and/or include the omitted
%   definitions in an appendix.
% Reason (claude): Lets reviewers check correctness claims and the timing statement; related PR-TECHNICAL-0011 and -0020.
% END-PAPER-REVIEW[codex:PR-STRUCTURE-0002,claude:PR-TECHNICAL-0023]
% PAPER-REVIEW[claude:PR-STRUCTURE-0003] pass=structure state=open
% Issue: The module format, the title's subject, is never specified; readers must infer it from examples.
% Evidence: The meaning of the doc/exports/imports/defs clauses, the "host" pseudo-module, the import-list form, and
%   the fact that defs evaluate to name/expression pairs are only implicit in Listing 1.1, Appendix C and the code
%   of eval-module. Listing 1.1 writes defs with "!fn" while Appendix C writes quoted pairs '(name ...), and the
%   connection (fn expands to a quoted pair) appears only in Listing C.2.
% Proposal: Open Section 4 with a compact description of the module format (one sentence per clause, how imports
%   are resolved, that each def must evaluate to a (name expr) pair, and what fn does), then present the loader.
% Reason: Gives the reader a specification to check the loader against, and connects Listing 1.1 with Appendix C.
% END-PAPER-REVIEW[claude:PR-STRUCTURE-0003]
We construct a program called \ild{module-loader} that is able to process
programs organised in modules like the one defined in
\cref{example-module-fact}. We will present it in a top-down manner while
omitting technical details such as \ild{let*} or \ild{kv-*} functions for
operating on key-value lists. The outermost expression defining the module
loader is shown in \cref{module-loader-outer}.
\begin{ildcode}
(!(eval (free-vars) (read-source
    "core/bootstrap/let.ild"))
  ( (host-bindings (free-vars))
    (lambda (eval (free-vars) (read-source
      "core/bootstrap/lambda-macro.ild")))
    (letrec (eval (free-vars) (read-source
      "core/bootstrap/letrec.ild")))
    (if (eval (free-vars) (read-source
      "core/bootstrap/if.ild")))
    (map (eval (free-vars) (read-source
      "core/bootstrap/map.ild"))))

  (!letrec (
    ; <definitions of helpers>
    ; <definition of module-loader>
    )

    module-loader))
\end{ildcode}
\lstcaption{Outer scope bindings of the module loader}\label{module-loader-outer}
% PAPER-REVIEW[codex:PR-TECHNICAL-0008,claude:PR-TECHNICAL-0011] pass=technical state=open
% Issue (codex): The displayed loader bootstrap does not bind let* although eval-module calls it.
% Evidence (codex): module-loader-outer binds lambda, letrec, if and map; let* is absent. eval-module uses !let*. The local core/bootstrap/module-loader.ild supplies a let* binding from let-star.ild.
% Proposal (codex): Include the let* bootstrap binding or explicitly include an equivalent definition among the omitted helpers and show how it reaches eval-module's captured environment.
% Reason (codex): The printed construction otherwise leaves an essential macro unresolved.
% Issue (claude): [Demonstrated inconsistency] Listings use bindings that are neither defined in the paper nor in scope.
% Evidence (claude): Listing 4.1 binds only host-bindings, lambda, letrec, if, map, but Listings 4.2-4.3 use let*, list,
%   cadr (via let), concat, sym-item-pair and a macro !and, which is never defined. The text says it omits "let*",
%   yet let* is now given in Listing 3.10. Listing 3.9 calls "add", but the host provides "+" (Section 2.3.1).
%   Listing C.1 imports gensym from host, which Section 2.3 does not provide; Section 7.1.1 says it needs an impure host.
% Proposal (claude): Add the missing bindings to Listing 4.1 (or say explicitly that it is abridged and list what is omitted,
%   including and/concat/sym-item-pair), replace add by +, and drop gensym from Listing C.1 or add it to Section 2.3.
% Reason (claude): The paper claims everything is built from the stated host environment; the listings should reflect that.
% END-PAPER-REVIEW[codex:PR-TECHNICAL-0008,claude:PR-TECHNICAL-0011]
The core of the module loader is a function called \ild{eval-module} (\cref{def-eval-module})
that:
\begin{enumerate}
\item Recursively processes imported modules and builds a binding environment \ild{env}
that contains all imports
\item Builds a list of name-body pairs from the definitions \paraphrase{given} in the module
\item Calls \ild{letrec} with these definition pairs and a body that generates a list
of key-value pairs for the module exports. Note that unlike other places, where letrec
is used as a macro, here it is used as a function (\cref{macro-function-duality}).
\item Evaluates the resulting program under \ild{env}
\end{enumerate}
\begin{ildcode}
; ---- inside letrec block ----
(load-module (!lambda (module-name)
  (eval-module (read-source module-name))))

(eval-module (!lambda (mod-expr)
  (!let* (
    (items (cdr mod-expr))

    (imports (kv-fetch 'imports (list) items))
    (defs    (kv-fetch 'defs    (list) items))
    (exports (kv-fetch 'exports (list) items))

    (env (process-imports imports))
    (eval-def (!lambda (def)
      (eval env def)))
    (def-items (map eval-def defs)))

    (eval env
      (letrec
        def-items
        (cons list
          (map sym-item-pair exports)))))))
\end{ildcode}
\lstcaption{Definitions of \ild{load-module} and \ild{eval-module}}\label{def-eval-module}
% PAPER-REVIEW[codex:PR-TECHNICAL-0014] pass=technical state=open
% Issue: The loader's domain and the assumptions behind its purity/caching claims are underspecified.
% Evidence: process-imports concatenates bindings although Env assumes unique keys; load-module recursively follows imports without a cycle case. read-source is described as filesystem IO, while future-work calls the host purely functional and caching only an efficiency concern.
% Proposal: State a finite acyclic import graph and unique-import-name precondition or define conflict/cycle handling; model read-source as a fixed source mapping if purity is intended, and condition caching equivalence on the allowed host effects and module initialization.
% Reason: These assumptions determine whether the loader produces a valid environment, terminates, and preserves observable behaviour under caching.
% END-PAPER-REVIEW[codex:PR-TECHNICAL-0014]
The \ild{process-imports} helper (\cref{def-process-imports}) calls \ild{load-module} recursively on
imported modules and then merges the resulting export lists. For the special
case when the \ild{host} keyword is used in place of a module name,
the host bindings (obtained via an outermost call to \ild{free-vars}) are
used in place of an export list.
\begin{ildcode}
; ---- inside letrec block ----
(process-imports (!lambda (import-list)
  (concat (map process-import import-list))))

(process-import (!lambda (import-spec)
  (!if (!and (symbol? (car import-spec))
             (sym-eq? (car import-spec) 'host))
    (process-import-items
      host-bindings
      (cdr import-spec))
    (process-import-items
      (load-module (car import-spec))
      (cdr import-spec)))))
\end{ildcode}
\lstcaption{Definition of \ild{process-imports}}\label{def-process-imports}
We omit the implementation of \ild{process-import-items}, which merely
picks the requested key-value pairs from the export lists.

Finally, we implement the function \ild{module-loader} which, given a module name,
entrypoint function name and arguments, loads that module, picks the entrypoint function
from its export list and applies it to the given arguments (implementation in \cref{def-module-loader}).
\begin{ildcode}
; ---- inside letrec block ----
(module-loader
  (!lambda
    (mod-name entry-func-name . args)
    (apply
      (kv-get entry-func-name
        (load-module mod-name))
      args)))
\end{ildcode}
\lstcaption{Definition of \ild{module-loader}}\label{def-module-loader}
Running a program
such as the one in \cref{example-module-fact} is now a matter of passing its
name and entrypoint to the module loader (shown in \cref{module-load-example}).
\begin{ildcode}
((eval (free-vars) (read-source
  "core/bootstrap/module-loader.ild"))

  "demos/fact.ild"
  'main)
\end{ildcode}
\lstcaption{Executing the module loader on \paraphrase{a user} program}\label{module-load-example}
The ILD interpreter and/or build system need not be \paraphrase{concerned with} any details
about loading and importing modules and can be kept simple: all of the logic
is implemented in ILD itself.
